\begin{Abstract}
In exact arithmetic, the expression $(x^{10}+c)-x^{10}$ is identically equal to $c$. In finite-precision binary arithmetic, however, its computed value displays a sequence of plateaus, parity-dependent rounding ties, and eventually complete absorption of the increment. This report investigates that behavior in IEEE 754 \code{binary32} and \code{binary64}, using the supplied research manuscript, extended book draft, experimental report, notebook, and interactive graphics. Writing $y=\fl(x^{10})$, $u=\ulp(y)$, and $c=m2^{\nu}$ for odd $m$, the nearest-grid model predicts an initial inexact binade at $E_1=p+\nu$ and the last potentially nonzero binade at $E_2=p+\lfloor\log_2 c\rfloor$, where $p$ is significand precision. Consequently, the predicted inexact interval in $x$ is approximately $[2^{(p+\nu)/10},\,2^{(p+1+\lfloor\log_2 c\rfloor)/10})$ and spans $L=\lfloor\log_2 c\rfloor-\nu+1$ binades. The model is established under explicit assumptions about rounding, reuse of the computed power, and binade boundaries. Results reported in the supplied study cover 1,611 values of $c$ in each precision and distinguish powers of two, odd integers, dyadic fractions, and non-dyadic decimal or reciprocal inputs. Reconstructed visualizations and an independently checked 780-case power-of-two grid illustrate when increments are preserved, rounded, or absorbed. Special attention is given to boundary effects, input representation, the relation to Fast2Sum, and limitations of the model.

\Keywords IEEE 754; floating-point arithmetic; unit in the last place; rounding to nearest; cancellation; absorption; binary32; binary64; numerical reproducibility.
\end{Abstract}

\tableofcontents
\listoffigures
\listoftables
